\QHeading{第1题　口算。}
\par\noindent


\begin{oralgrid}
\OralItem{1}{4\times 25}{100}
\OralItem{2}{25\times 8}{200}
\OralItem{3}{8\times 125}{1000}
\OralItem{4}{125\times 4}{500}
\OralItem{5}{50\div 2}{25}
\OralItem{6}{2\times 50}{100}
\OralItem{7}{24\times 4}{96}
\OralItem{8}{50\times 24}{1200}
\OralItem{9}{560\div 8}{70}
\OralItem{10}{360\div 9}{40}
\OralItem{11}{35\times 20}{700}
\OralItem{12}{40\times 45}{1800}\end{oralgrid}
\par

\QHeading{口算思路补充（对应上表题号）}

\Answer{上面列出得数，下面按小题给出不同计算思路。计数单位法适用于小数，整数题可比较凑整、拆分或商不变等方法。}
\MethodTitle{第1小题：$4\times 25$}

\MethodBlock{口算思路一：按数位拆分计算}{把一个因数拆成整十、整百数与余下部分之和。}{\begin{aligned}[t]&\phantom{={}}4\times 25\\
&=4\times (20+5)\\
&=4\times 20+4\times 5\\
&=80+20\\
&=100\end{aligned}}
\MethodTitle{第2小题：$25\times 8$}

\MethodBlock{口算思路一：拆因数，凑整计算}{把一个因数拆成两个因数，先凑成整百或整千。}{\begin{aligned}[t]&\phantom{={}}25\times 8\\
&=25\times (4\times 2)\\
&=(25\times 4)\times 2\\
&=100\times 2\\
&=200\end{aligned}}

\MethodBlock{口算思路二：凑整后补偿}{把接近整十、整百的因数凑整，再补回或减去多算的部分。}{\begin{aligned}[t]&\phantom{={}}25\times 8\\
&=25\times (10-2)\\
&=25\times 10-25\times 2\\
&=250-50\\
&=200\end{aligned}}
\MethodTitle{第3小题：$8\times 125$}

\MethodBlock{口算思路一：凑整后补偿}{把接近整十、整百的因数凑整，再补回或减去多算的部分。}{\begin{aligned}[t]&\phantom{={}}8\times 125\\
&=125\times (10-2)\\
&=125\times 10-125\times 2\\
&=1250-250\\
&=1000\end{aligned}}

\MethodBlock{口算思路二：按数位拆分计算}{把一个因数拆成整十、整百数与余下部分之和。}{\begin{aligned}[t]&\phantom{={}}8\times 125\\
&=8\times (100+25)\\
&=8\times 100+8\times 25\\
&=800+200\\
&=1000\end{aligned}}
\MethodTitle{第4小题：$125\times 4$}

\MethodBlock{口算思路一：按数位拆分计算}{把一个因数拆成整十、整百数与余下部分之和。}{\begin{aligned}[t]&\phantom{={}}125\times 4\\
&=4\times (100+25)\\
&=4\times 100+4\times 25\\
&=400+100\\
&=500\end{aligned}}
\MethodTitle{第5小题：$50\div 2$}

\MethodBlock{口算思路一：拆分被除数}{}{\begin{aligned}[t]&\phantom{={}}50\div 2\\
&=(40+10)\div 2\\
&=40\div 2+10\div 2\\
&=20+5\\
&=25\end{aligned}}

\MethodBlock{口算思路二：利用商不变的性质}{被除数和除数同时乘同一个非零数，商不变。}{\begin{aligned}[t]&\phantom{={}}50\div 2\\
&=(50\times 5)\div (2\times 5)\\
&=250\div 10\\
&=25\end{aligned}}
\MethodTitle{第7小题：$24\times 4$}

\MethodBlock{口算思路一：按数位拆分计算}{把一个因数拆成整十、整百数与余下部分之和。}{\begin{aligned}[t]&\phantom{={}}24\times 4\\
&=4\times (20+4)\\
&=4\times 20+4\times 4\\
&=80+16\\
&=96\end{aligned}}
\MethodTitle{第8小题：$50\times 24$}

\MethodBlock{口算思路一：拆因数，凑整计算}{把一个因数拆成两个因数，先凑成整百或整千。}{\begin{aligned}[t]&\phantom{={}}50\times 24\\
&=50\times (2\times 12)\\
&=(50\times 2)\times 12\\
&=100\times 12\\
&=1200\end{aligned}}

\MethodBlock{口算思路二：按数位拆分计算}{把一个因数拆成整十、整百数与余下部分之和。}{\begin{aligned}[t]&\phantom{={}}50\times 24\\
&=50\times (20+4)\\
&=50\times 20+50\times 4\\
&=1000+200\\
&=1200\end{aligned}}
\MethodTitle{第9小题：$560\div 8$}

\MethodBlock{口算思路一：利用商不变的性质}{被除数和除数同时乘同一个非零数，商不变。}{\begin{aligned}[t]&\phantom{={}}560\div 8\\
&=(560\times 5)\div (8\times 5)\\
&=2800\div 40\\
&=70\end{aligned}}

\MethodBlock{口算思路二：拆分除数，连续相除}{}{\begin{aligned}[t]&\phantom{={}}560\div 8\\
&=560\div (2\times 4)\\
&=560\div 2\div 4\\
&=280\div 4\\
&=70\end{aligned}}
\MethodTitle{第10小题：$360\div 9$}

\MethodBlock{口算思路一：拆分除数，连续相除}{}{\begin{aligned}[t]&\phantom{={}}360\div 9\\
&=360\div (3\times 3)\\
&=360\div 3\div 3\\
&=120\div 3\\
&=40\end{aligned}}
\MethodTitle{第11小题：$35\times 20$}

\MethodBlock{口算思路一：按数位拆分计算}{把一个因数拆成整十、整百数与余下部分之和。}{\begin{aligned}[t]&\phantom{={}}35\times 20\\
&=20\times (30+5)\\
&=20\times 30+20\times 5\\
&=600+100\\
&=700\end{aligned}}
\MethodTitle{第12小题：$40\times 45$}

\MethodBlock{口算思路一：按数位拆分计算}{把一个因数拆成整十、整百数与余下部分之和。}{\begin{aligned}[t]&\phantom{={}}40\times 45\\
&=40\times (40+5)\\
&=40\times 40+40\times 5\\
&=1600+200\\
&=1800\end{aligned}}


\par\noindent

\QHeading{第2题　在$\square$里填上合适的数或字母。}

\Equation{36\times 5\times 2=\textcolor{answerred}{36}\times (\textcolor{answerred}{5}\times 2)}
\par

\Equation{(25\times 12)\times 4=(25\times\textcolor{answerred}{4})\times\textcolor{answerred}{12}}
\Equation{(a\times b)\times c=\textcolor{answerred}{a}\times (b\times\textcolor{answerred}{c})}
\Equation{125\times 23\times 8=(125\times\textcolor{answerred}{8})\times\textcolor{answerred}{23}}
\Equation{18\times 50\times 2=\textcolor{answerred}{18}\times (\textcolor{answerred}{50}\times 2)}
\Equation{(82\times 4)\times 250=\textcolor{answerred}{82}\times (\textcolor{answerred}{4}\times\textcolor{answerred}{250})}
\par\noindent

\QHeading{第3题　用简便方法计算下面各题。}

\CalcBlock[方法一：简便计算]{1}{\begin{aligned}[t]&\phantom{={}}\color{black}18\times 25\times 40\\&\color{answerred}=18\times(25\times40)\\&\color{answerred}=18\times1000\\&\color{answerred}=18000\end{aligned}}
\par
\MethodBlock{方法二：结合因数，先凑整}{用乘法交换律、结合律，先算容易凑整的两个因数。}{\begin{aligned}[t]&\phantom{={}}18\times 25\times 40\\
&=(18\times 40)\times 25\\
&=720\times 25\\
&=18000\end{aligned}}

\MethodBlock{方法三：按运算顺序直接计算（对照）}{本法只作结果核对；本题要求简便计算，前面的简便解法仍保留。}{\begin{aligned}[t]&\phantom{={}}18\times 25\times 40\\
&=450\times 40\\
&=18000\end{aligned}}


\CalcBlock[方法一：简便计算]{2}{\begin{aligned}[t]&\phantom{={}}\color{black}125\times 4\times 8\times 25\\&\color{answerred}=(125\times8)\times(4\times25)\\&\color{answerred}=1000\times100\\&\color{answerred}=100000\end{aligned}}
\MethodBlock{方法二：结合因数，先凑整}{用乘法交换律、结合律，先算容易凑整的两个因数。}{\begin{aligned}[t]&\phantom{={}}125\times 4\times 8\times 25\\
&=(125\times 8)\times 4\times 25\\
&=1000\times 4\times 25\\
&=4000\times 25\\
&=100000\end{aligned}}

\MethodBlock{方法三：先处理括号或乘除，再拆因数，凑整计算}{把一个因数拆成两个因数，先凑成整百或整千。}{\begin{aligned}[t]&\phantom{={}}125\times 4\times 8\times 25\\
&=500\times 8\times 25\\
&=4000\times 25\\
&=25\times (4\times 1000)\\
&=(25\times 4)\times 1000\\
&=100\times 1000\\
&=100000\end{aligned}}

\MethodBlock{方法四：先处理括号或乘除，再把特殊因数改成整百或整千}{先把25、125等改写成整百或整千除以一个数。}{\begin{aligned}[t]&\phantom{={}}125\times 4\times 8\times 25\\
&=500\times 8\times 25\\
&=4000\times 25\\
&=(100\div 4)\times 4000\\
&=100\times (4000\div 4)\\
&=100\times 1000\\
&=100000\end{aligned}}

\CalcBlock[方法一：简便计算]{3}{\begin{aligned}[t]&\phantom{={}}\color{black}450\div 25\div 2\\&\color{answerred}=450\div(25\times2)\\&\color{answerred}=450\div50\\&\color{answerred}=9\end{aligned}}
\MethodBlock{方法二：先处理括号或乘除，再利用商不变的性质}{被除数和除数同时乘同一个非零数，商不变。}{\begin{aligned}[t]&\phantom{={}}450\div 25\div 2\\
&=18\div 2\\
&=(18\times 5)\div (2\times 5)\\
&=90\div 10\\
&=9\end{aligned}}

\CalcBlock[方法一：简便计算]{4}{\begin{aligned}[t]&\phantom{={}}\color{black}540\div (9\times 2)\\&\color{answerred}=540\div9\div2\\&\color{answerred}=60\div2\\&\color{answerred}=30\end{aligned}}
\MethodBlock{方法二：先处理括号或乘除，再拆分除数，连续相除}{}{\begin{aligned}[t]&\phantom{={}}540\div (9\times 2)\\
&=540\div 18\\
&=540\div (2\times 9)\\
&=540\div 2\div 9\\
&=270\div 9\\
&=30\end{aligned}}

\MethodBlock{方法三：先处理括号或乘除，再利用商不变的性质}{被除数和除数同时乘同一个非零数，商不变。}{\begin{aligned}[t]&\phantom{={}}540\div (9\times 2)\\
&=540\div 18\\
&=(540\times 5)\div (18\times 5)\\
&=2700\div 90\\
&=30\end{aligned}}

\CalcBlock[方法一：简便计算]{5}{\begin{aligned}[t]&\phantom{={}}\color{black}1000\div (25\times 5)\\&\color{answerred}=1000\div125\\&\color{answerred}=8\end{aligned}}
\MethodBlock{方法二：除以积，改成连续除}{连续除以两个数，等于除以这两个数的积。}{\begin{aligned}[t]&\phantom{={}}1000\div (25\times 5)\\
&=1000\div 25\div 5\\
&=40\div 5\\
&=8\end{aligned}}

\CalcBlock[方法一：简便计算]{6}{\begin{aligned}[t]&\phantom{={}}\color{black}800\div 16\div 5\\&\color{answerred}=800\div(16\times5)\\&\color{answerred}=800\div80\\&\color{answerred}=10\end{aligned}}
\MethodBlock{方法二：先处理括号或乘除，再利用商不变的性质}{被除数和除数同时乘同一个非零数，商不变。}{\begin{aligned}[t]&\phantom{={}}800\div 16\div 5\\
&=50\div 5\\
&=(50\times 2)\div (5\times 2)\\
&=100\div 10\\
&=10\end{aligned}}

\par\noindent

\QHeading{第4题　解决问题。}

\Problem{果园里有2600千克香蕉需要装箱运走，平均每箱装25千克。一辆小推车一次能运4箱，这些香蕉需要用这辆小推车运多少次?}
\par


\MethodTitle{方法一}
\Answer{\Formula{\begin{aligned}[t]&\phantom{={}}2600\div(25\times4)\\&=2600\div100\\&=26\text{（次）}\end{aligned}}}
\Answer{答：需要运26次。}

\MethodTitle{方法二：先求箱数，再求运输次数}
\Answer{先按每箱25千克装箱，再按一次4箱计算。}
\Answer{\Formula{\begin{aligned}[t]&\phantom{={}}2600\div 25\div 4\\
&=104\div 4\\
&=26\text{（次）}\end{aligned}}}
\Answer{答：需要运26次。}
